\begin{remark}[where the quadratic-like embedding becomes only an immersion]\label{rem:kahn-immersion-disclosures}
The sources disclose, at four separate places, that the
inclusion-type arrow of the quadratic-like theory degenerates to an
immersion in the \(\psi\)-quadratic-like theory; none of these places
asserts a global embedding of the localized inclusion.
\begin{enumerate}
\item \textbf{Kahn~I, \S6.1, pp.~31--32.}  In the
quadratic-like sections the map
\(i\colon U^{n+1}\sm\KK\to U^n\sm\KK\) is literally the inclusion;
footnote~15 on p.~31 announces that in the \(\psi\)-quadratic-like
sections (\S\S7--8) this embedding will be
replaced by an immersion, ``so it is good to have notation for it.''
The same passage notes that the triangle
formed by the covering \(f\), the inclusion
\(U^{n+1}\sm\widetilde\KK\subset U^{n+1}\sm\KK\), and \(i\) does
\emph{not} commute; passage along that hole-filling leg is controlled
quantitatively via Property~D (Kahn~I, \S2.7).
\item \textbf{Kahn~I, \S1.4.}
\(\psi\)-quadratic-like maps are introduced as the generalization in
which the inclusion \(\BU\to\BV\) is replaced by an immersion \(i\)
whose properties imply only that the multivalued \(f\circ i^{-1}\) has
a single-valued quadratic-like \emph{restriction}.
\item \textbf{Kahn~I, \S8.1.}  At each stage of
the \(\psi\)-quadratic-like tower there are three maps: a covering
\(f\colon\BU^{n+1}\sm f^{-1}\KK\to\BU^n\sm\KK\); an inclusion
\(\BU^{n+1}\sm f^{-1}\KK\hookrightarrow\BU^{n+1}\sm\KK\); and an
immersion \(i\colon\BU^{n+1}\sm\KK\to\BU^n\sm\KK\).  The source
states that the one significant change from the quadratic-like setting
is that the third map, previously an embedding, is now an immersion.
The subsequent proofs use, about \(i\), only (a)
properness on the ends corresponding to \(\KK\) and (b)
that the iterated immersion is homotopic to a homeomorphism;
point-injectivity of \(i\) is neither claimed nor used.
\item \textbf{KL~II, \S3.2.2, Lemma~3.4, pp.~12--13.}  The naturality lemma
keeps two arrows strictly separate: the dynamical arrow
\(\mathbf f\colon\mathbf Y\to\mathbf Z\) induced by
\(f\colon Y\to Z\) is a branched covering of the same degree \(k\) (an
isomorphism of glued universal-cover copies when \(k=1\);
Lemma~3.4(i), p.~13), while the inclusion arrow
\(\mathbf i\colon\mathbf Y\to\mathbf Z\) induced by \(Y\subset Z\) is
an immersion built from component-wise immersions of universal covers,
with the component assignment \(j(i)\) described as ``surjective but
not necessarily injective'' (proof of Lemma~3.4(ii), p.~13; p.~68 of
the journal pagination).  This is the closest published anticipation of
Theorem~\ref{thm:forced-collision}, and it is discussed as such in
Section~\ref{subsec:literature}, where the three ways in which the
present theorem differs from it are set out.  Ideal-boundary behavior
is likewise not embedding-like: the boundary extension of
\(\mathbf i\) exists, but KL state that
\(\mathbf i(\partial\mathbf Y)\) is not contained in
\(\partial\mathbf Z\) unless \(\mathbf Z=\mathbf Y\)
(KL~II, \S3.2.4, p.~13) --- a stronger assertion than the ``need not''
used elsewhere in this paper --- and the extremal-length
arguments of KL~II use maximal initial lifts that may terminate on the
\(O\)-boundary (KL~II, \S5.5).
\end{enumerate}
\end{remark}

\begin{proof}[Source verification (documentary; not a mathematical proof)]
Each item paraphrases the cited passage of the published sources:
items (1)--(3) against Kahn~I, arXiv:math/0609045v3 (28 December 2023);
item (4) against KL~II, Ann.\ Sci.\ \'Ec.\ Norm.\ Sup\'er.\ (4)
\textbf{41} (2008) 57--84, and the identically sectioned
arXiv:math/0609046v2.  No further claim is made: in
particular none of the cited passages asserts, and no other passage of
either source asserts, that the localized inclusion
\(\widehat i\) of dynamical localization is globally injective.  The
five-word phrase quoted in item (4) is the only verbatim material
retained from either source.
\end{proof}
